\uuid{xvN1}
\exo7id{4481}
\titre{exo7 4481}
\auteur{quercia}
\organisation{exo7}
\datecreate{2010-03-14}
\isIndication{false}
\isCorrection{true}
\chapitre{Série numérique}
\sousChapitre{Autre}
\cpge{mpsi-mp2i.nombres-reels-suites.d, pcsi.nombres-reels-suites.c, ptsi.nombres-reels-suites.c, mpsi-mp2i.procedes-sommatoires-discrets.a, pcsi.series-numeriques.a, ptsi.series-numeriques.a, bcpst1.suites-reelles, mp-mpi.series-numeriques-vectorielles.a, pt.series-numeriques.b, mp-mpi.series-numeriques-vectorielles.b, pc.series-numeriques, psi.series-numeriques.a, pt.series-numeriques.a, bcpst2.series-reelles}
\module{Analyse}
\niveau{L1}
\difficulte{}

\contenu{
\question{
Soit $(u_n)$ une suite de complexes telle que $\frac{u_1+\dots+u_n}n\to\ell\in\C$ lorsque $n\to\infty$.
Montrer que $\frac{1}{\ln(n)}\Bigl(\frac{u_1}1+\dots+\frac{u_n}n\Bigr)\to \ell$ lorsque $n\to\infty$.
}
\reponse{
Transformaton d'Abel.
}
}
